\documentclass[../../main.tex]{subfiles}

\begin{document}

\section{Attention and Transformers}

Finally! Its time to discuss about transformers. If you have not
already, please complete the playlist by 3Blue1Brown's highlighted
in the last note and read \href{https://arxiv.org/abs/1706.03762}
{Attention Is All You Need}. More great resources can be found in
the last note. \textit{Attention Is All You Need} is considered one
of the most important paper in computer science and behind its
meme-like title, it makes current LLMs possible.

Although the transformer architecture is used in many places, I want
to introduce it LLMs as I think it is most intuitive and classic.
The problem with prior language models is that they can't remember
what they were talking about seconds ago. Because the next token
is predicted with limited context window, it is not possible to
``remember'' the entire context like us. To improve this, eight
researchers from Google introduced transformers. Although the basic
ideas existed before the paper, it made the transformers well known.

The content itself, when you think about it, is quite simple and
very intuitive. What should we do to enhance the use of information
that we have? The answer is quite simple and circular. Literally use
more prior contexts with efficient methods. Think of transformer
blocks, or layers, as a function that takes in and spits out vectors
of the same shape. The input vector goes in to a magical concept
called multi-head attention, add and norm, and feed forward (which is
basically a fancier word for MLPs that we discussed so far).


\end{document}


